\documentclass[aps,prb,twocolumn,superscriptaddress,reprint]{revtex4-2}
\usepackage{amsmath,amssymb}
\usepackage{graphicx}
\begin{document}
\title{Anomalous Thermal Hall Response in Twisted Bilayer Magnonic Crystals}
\author{Mirei Kasahara}
\affiliation{Institute for Imaginary Condensed Matter, Kitamachi University, Kitamachi 100-0000, Japan}
\author{Daniel R. Oakfield}
\affiliation{Department of Physics, Westbrook Institute of Technology, Westbrook 00000, Freedonia}
\author{Lena Varga}
\affiliation{Institute for Imaginary Condensed Matter, Kitamachi University, Kitamachi 100-0000, Japan}
\affiliation{Center for Quantum Fictions, Westbrook 00000, Freedonia}
\date{\today}
\begin{abstract}
We study the thermal Hall effect of magnons in twisted bilayers of a fictional honeycomb ferromagnet.
Using a linear spin-wave theory combined with a continuum model of the moir\'e superlattice, we show that the Berry curvature of the lowest magnon bands becomes strongly concentrated near small twist angles.
The resulting thermal Hall conductivity changes sign as a function of temperature and reaches values an order of magnitude larger than in the untwisted bilayer.
We identify the origin of this enhancement in a cascade of band inversions and propose an experimental protocol based on local thermometry.
Our results suggest that twist engineering offers a simple route to tune topological magnon transport.
\end{abstract}
\maketitle
\section{Introduction}
Topological magnons have attracted growing attention as carriers of heat that are free from Joule dissipation~\cite{kasahara2021,oakfield2022}.
In layered magnets, the relative twist between layers provides an additional knob that reshapes the magnon band structure without changing the chemistry of the material.
Previous studies focused on the magnon spectrum itself, leaving the transport consequences largely unexplored~\cite{varga2023}.

In this paper we fill this gap by computing the thermal Hall conductivity of twisted bilayers over a wide range of twist angles and temperatures.
We find that the response is dominated by a small number of flat bands whose Berry curvature can be enhanced by more than a factor of ten.
The mechanism is generic and should apply to a broad class of van der Waals magnets, including several candidates that have been synthesized recently.
We also discuss how disorder and magnon--phonon coupling modify these conclusions and estimate the size of the signal in realistic devices.

\section{Model}
The bilayer is described by a Heisenberg Hamiltonian with Dzyaloshinskii--Moriya interactions.
\bibliography{refs}
\end{document}
